\scanpage{27}\mspage{26}
% L1–5
Les tronçons $]a,b[,\ L_{<a},\ L_{>a},\ \inserted{L}$ sont dits ouverts \deleted{(\uline{NB} ils sont} \deleted{(ce sont bien ceux} \inserted{(ce sont bien ceux qui sont} ouverts pour la top. can.), \deleted{qui sont ouverts} et $]a,b]$ et $[b,\unsure{o|a|0}[$ sont dits \uline{demi}-\uline{ouverts}.

% L6–11
\uline{Attention} \deleted{Le tronçon} Si $a$ plus petit \unsure{élém|élément}, alors $[a,b[$, \deleted{est ouvert} \inserted{qui est} demi-ouvert pour la terminologie \unsure{précédente}, est \unsure{aussi} ouvert pour la top. canonique, et il est égal à $L_{<b}$. Kif kif \unsure{dualement|autrement}. $L$ est à la fois ouvert et fermé.

% L12–14
Un \uline{tronçon} \uline{générali\deleted{\unsure{sé}}\unsure{sé}} \inserted{ou \unsure{ps.tronçon}} est une partie de $L$ qui est soit \unsure{réduit|réduite} à un pt, soit un tronçon. Pour un tronçon de la forme
% L15–16
\[
\deleted{\unsure{T}}\ T'=T\setminus\alpha\qquad(T\text{ tronçon fermé},\ \alpha\in\partial T
\]
je pose
\[
\partial T'=\partial T
\]
% L17–18
et \unsure{on aura|l'on aura} \unsure{encore|que}
\[
\partial T'=T'\mathbin{\unsure{\dot\cup|\dot\cap|\cup}}\partial L\deleted{\unsure{.}}
\]
\margin{Attention, cette fois ce n'est pas le $\partial$ absolu de $T'$ pour la structure d'ordre induite. \unsure{\ldots}}
% slanted
% L lower left, oblique
\margin{\uline{NB} $\partial T'\neq\varnothing$ sauf \unsure{qd|q} si $T'=L$, $L$ \unsure{une|un} \deleted{\unsure{``droite ouverte''}} \deleted{\unsure{ligne}} \unsure{tronçon-interminable}\ldots\ Les pts frontière de $T'$ sont \deleted{ses} bornes inf. et sup. \uline{dans $T$} (si elles existent) : borne sup (ssi $T'$ est majoré) borne inf (ssi $T'$ est minoré) égales \unsure{ds} le cas $T$ \unsure{réduit} à un pt\ldots}

% L19–22
\deleted{\unsure{Soient} deux} Deux \unsure{ps.tronçons} \inserted{$T,T'$} sont dits \uline{en pos. relative modérée} (\unsure{abrég}, $(T,T')$ \deleted{\unsure{en p.}} modéré ou $\{T\unsure{\mid|,}T'\}$ modéré) si $\partial T$ et $\partial T'$ le sont, i.e $\partial T\cup\partial T'$ est tot. ordonné.
